\section*{CHAPTER 3: IMPLEMENTATION AND DEPLOYMENT}
\addcontentsline{toc}{section}{\numberline{}CHAPTER 3: IMPLEMENTATION AND DEPLOYMENT}
\setcounter{section}{3}
\setcounter{subsection}{0}
\setcounter{figure}{0}
\setcounter{table}{0}

\subsection{Provider and Tenant Administration Modules}

The administration interfaces were built to prioritize usability and data density. The \textbf{Provider Admin} dashboard provides a high-level overview of the SaaS platform. It allows the super admin to monitor active tenants, configure subscription packages (setting limits on the number of fields or staff accounts per package), and track monthly revenue generated from these subscriptions.

The \textbf{Tenant Admin} dashboard serves as the operational hub for facility owners. A key feature of this interface is the \textbf{Time-grid Schedule Viewer}. This component provides a visual representation of all football fields over the course of a day. The grid is dynamically generated based on existing bookings in the database, using color-coded blocks to indicate available, pending, confirmed, and maintenance statuses. This visual approach effectively eliminates the risk of double-booking a field.

Furthermore, the Tenant module includes comprehensive CRM capabilities, automatically classifying customers into "VIP" or "New" based on their total spending. The financial module calculates occupancy rates and generates printable invoices for completed bookings.

\subsection{Customer Storefront and AI Integration}

The \textbf{Customer Storefront} is a responsive web application designed for mobile-first usage. It dynamically loads the theme, logo, and available fields based on the specific \texttt{tenant\_id} present in the URL, allowing each facility to have a personalized storefront without requiring separate codebases.

The most innovative feature of the storefront is the \textbf{AI Chatbot}, powered by Google Gemini 2.5 Flash. The chatbot interface floats seamlessly over the web page. When a user asks a question, such as "Are there any 7-a-side fields available tomorrow at 18:00?", the AI translates this intent into a structured "Function Call". The Node.js backend intercepts this call, queries the database, and returns the result to the AI, which then replies: "Yes, Field B is available. Would you like me to book it for you?". If the user confirms, the AI triggers the \texttt{create\_booking} function, finalizing the transaction.

\subsection{Payment Integration and Mobile Application}

To secure bookings, the platform integrates the \textbf{VNPay} payment gateway. Upon booking confirmation, users are redirected to VNPay to complete the transaction. A secure callback URL processes the payment status and updates the database accordingly, issuing an electronic QR ticket to the user.

To bridge the gap between web software and physical operations, the web application is packaged into a native Android application using \textbf{Capacitor}. This allows the Tenant Admin app to access the device's native camera. Facility staff can simply open the app, tap the "Check-in QR" button, and scan the customer's electronic ticket to instantly verify and update the booking status in the database.

\subsection{Cloud Deployment}

The deployment pipeline is fully automated. The Node.js backend and static frontend files are hosted on \textbf{Render}, connected directly to the GitHub repository for Continuous Deployment (CD). The MySQL database is hosted on \textbf{Aiven}, providing a managed cloud environment with automatic backups and high availability, ensuring the SaaS platform remains stable under load.